\documentclass[12pt,aspectratio=169]{beamer}

\usetheme{Madrid}
\usecolortheme{default}
\usefonttheme{professionalfonts}

\usepackage[T2A]{fontenc}
\usepackage[utf8]{inputenc}
\usepackage[russian,english]{babel}
\usepackage{amsmath,amssymb,amsthm,mathtools}
\usepackage{array,booktabs}
\usepackage{tikz}
\usepackage{tikz-cd}
\usepackage{csquotes}
\usepackage{microtype}

\usepackage[
  backend=biber,
  style=authoryear,
  sorting=nyt,
  maxbibnames=4,
  giveninits=true,
  uniquename=init,
  uniquelist=false,
  doi=false,
  isbn=false,
  url=false,
  eprint=false,
  autolang=other
]{biblatex}
\addbibresource{../disser.bib}
\addbibresource{../disser-my.bib}

\definecolor{MainBlue}{RGB}{28,63,112}
\definecolor{MainRed}{RGB}{150,40,40}

\setbeamercolor{structure}{fg=MainBlue}
\setbeamercolor{title}{fg=MainBlue}
\setbeamercolor{frametitle}{fg=MainBlue,bg=MainBlue!8}
\setbeamercolor{block title}{fg=MainBlue,bg=MainBlue!12}
\setbeamercolor{block body}{fg=black,bg=black!2}
\setbeamercolor{exampleblock title}{fg=black,bg=black!10}
\setbeamercolor{exampleblock body}{fg=black,bg=black!1}
\setbeamercolor{alerted text}{fg=MainRed}

\setbeamertemplate{navigation symbols}{}
\setbeamertemplate{footline}[frame number]
\setbeamersize{text margin left=5mm,text margin right=5mm}

\setbeamerfont{title}{size=\Large}
\setbeamerfont{subtitle}{size=\normalsize}
\setbeamerfont{author}{size=\normalsize}
\setbeamerfont{institute}{size=\small}
\setbeamerfont{date}{size=\small}
\setbeamerfont{frametitle}{size=\large}
\setbeamerfont{block title}{size*={12pt}{14pt}}
\setbeamerfont{block body}{size=\normalsize}
\setbeamerfont{block title alerted}{size*={12pt}{14pt}}
\setbeamerfont{block body alerted}{size=\normalsize}
\setbeamerfont{block title example}{size*={12pt}{14pt}}
\setbeamerfont{block body example}{size=\normalsize}
\setbeamerfont{normal text}{size=\normalsize}

\DeclareMathOperator{\Ann}{Ann}
\DeclareMathOperator{\Aut}{Aut}
\DeclareMathOperator{\Der}{Der}
\DeclareMathOperator{\Char}{char}
\DeclareMathOperator*{\Supp}{Supp}

\newcommand{\K}{\mathbb{K}}
\newcommand{\quot}[2]{#1\mathbin{/}#2}
\newcommand{\term}[1]{{\color{MainRed}\bfseries #1}}
\DeclareCiteCommand{\slidecite}
  {}
  {%
    \printnames{author}%
    \setunit{\addcomma\space}%
    \iffieldundef{shortjournal}
      {%
        \iffieldundef{journaltitle}
          {%
            \iffieldundef{booktitle}
              {%
                \iffieldundef{publisher}
                  {\printfield{title}}
                  {\printlist{publisher}}%
              }
              {\printfield{booktitle}}%
          }
          {\printfield{journaltitle}}%
      }
      {\printfield{shortjournal}}%
    \setunit{\addcomma\space}%
    \printfield{year}%
  }
  {\addsemicolon\space}
  {}
\newcommand{\sourcefootnote}[1]{\footnote{\slidecite{#1}}}
\newcommand{\thmline}[1]{%
  \medskip
  {\color{MainBlue}\rule{\linewidth}{0.4pt}}
  \smallskip
  {\bfseries #1\par}
}

\AtBeginBibliography{\tiny}

\title[Левосимметрические структуры на $N\Phi(\K)$]{Левосимметрические структуры\\ на нильтреугольных алгебрах Шевалле}
\author[Н. Д. Ходюня]{Н. Д. Ходюня}
\institute[СФУ]{Сибирский Федеральный Университет}
\date[06.05.2026]{Семинар теории колец им. А.\,И. Ширшова\\6 мая 2026}

\begin{document}

\begin{frame}[plain]
    \vfill
    \begin{center}
        {\usebeamerfont{title}\usebeamercolor[fg]{title}\inserttitle\par}
        \vspace{1.6em}
        {\usebeamerfont{author}\insertauthor\par}
        \vspace{0.4em}
        {\usebeamerfont{institute}\insertinstitute\par}
    \end{center}
    \vfill
    \begin{center}
        {\usebeamerfont{date}\insertdate\par}
    \end{center}
\end{frame}

\begin{frame}[t]{План доклада}
    \begin{enumerate}
        \item Точные обертывающие алгебры для $N\Phi(\K)$ и постановка задачи.
        \item Градуированные алгебры и теорема единственности для типа $A_n$.
        \item Левосимметрические структуры, их полупрямые суммы и конструкции для типов
              $B_n$, $C_n$ и $D_n$.
    \end{enumerate}

    \medskip

    \begin{alertblock}{Основные результаты}
        \begin{itemize}
            %
            \item Если ассоциативная точная обертывающая алгебра $R$ алгебры Ли $N\Phi(\K)$ типа
                  $A_n$ $(n\ge 3)$ допускает градуировку, согласованную с градуировкой Картана,
                  то $R \simeq NT(n+1,\K)$.

            \item Для $N\Phi(\K)$ типов $B_n$, $C_n$ и $D_n$ существуют левосимметрические точные
                  обертывающие алгебры, построенные как полупрямые суммы алгебры
                  $NT(n,\mathbb{K})$ и подходящего идеала алгебры $N\Phi(\mathbb{K})$.
        \end{itemize}
    \end{alertblock}
\end{frame}

\begin{frame}[plain]
    \vfill
    \begin{center}
        {\usebeamerfont{title}\usebeamercolor[fg]{title}Точные обертывающие алгебры для $N\Phi(\K)$\\ и постановка задачи.\par}
    \end{center}
    \vfill
\end{frame}

\begin{frame}[t]{Неассоциативные точные обертывающие алгебры}
    \small
    Ассоциативное кольцо $R$ всегда превращается в кольцо Ли $R^{(-)}$, если
    заменить умножение в $R$ на коммутатор $[a,b]\coloneq ab - ba$.

    \begin{itemize}
        \item Произвольная (не обязательно ассоциативная) алгебра $R$ называется
              \term{Ли-допустимой}\sourcefootnote{albert48}, когда $R^{(-)}$ есть алгебра Ли.

        \item В.\,М.~Левчук\sourcefootnote{DAN2018,imm2020} называет алгебру $R$ \term{точной
                  обертывающей алгеброй} для алгебры Ли $L$, если
              \[
                  R^{(-)}\simeq L.
              \]
    \end{itemize}
\end{frame}

\begin{frame}[t]{Определение умножения в $R$ по заданной структуре $R^{(-)}$}
    \small
    Если $2\in \K^*$, то для любой $\K$-алгебры $R$ с умножением $xy$ имеем соотношение
    \[
        xy = \frac{1}{2}[x,y] + x \circ y, \text{ где } x \circ y \coloneq \frac{1}{2}(xy + yx).
    \]

    \begin{itemize}
        \item При заданной структуре на $R^{(-)}$ определение умножения $xy$ в $R$
              равносильно определению коммутативного произведения $x\circ y$ на векторном
              пространстве $R$.

        \item Поэтому при изучении Ли-допустимых алгебр основная задача состоит в определении
              коммутативных произведений $\circ$ по заданной структуре $R^{(-)}$ и
              дополнительным условиям на $R$.

        \item Вопрос об условиях однозначности точной обертывающей алгебры отмечал
              И.\,П.~Шестаков на конференции в 2017 году\sourcefootnote{modernWorldConf2017}.
    \end{itemize}
\end{frame}

\begin{frame}[t]{Нильтреугольные алгебры Шевалле}
    \small
    \begin{itemize}
        \item В теории Картана--Киллинга каждой простой комплексной конечномерной алгебре Ли
              $L$ сопоставляют единственную с точностью до эквивалентности неразложимую
              систему корней $\Phi$; выбор простых корней $\Pi$ определяет систему
              положительных корней $\Phi^+\supseteq \Pi$\sourcefootnote{carter72}.
        \item Элементы $e_r$ $(r \in \Phi)$ и подходящая база подалгебры Картана алгебры Ли
              $L=\mathfrak{L}(\Phi,\mathbb{C})$ задают базис Шевалле с целочисленными
              структурными константами\sourcefootnote{Chev-55}.
        \item Тем самым определяется алгебра Шевалле $\mathfrak{L}(\Phi,\K)$ над произвольным
              полем $\K$.
        \item Ее подалгебру $N\Phi(\K)$ с базой $\{e_r\mid r\in\Phi^+\}$ называем
              \term{нильтреугольной}.
    \end{itemize}
\end{frame}

\begin{frame}[t]{Универсальная и точная обертывающие}
    \small
    \begin{itemize}
        \item Согласно теореме Пуанкаре--Биркгофа--Витта каждая алгебра Ли может быть
              получена при помощи выбора подходящей ассоциативной обертывающей алгебры,
              замены в ней операции на коммутирование и дальнейшего перехода к подходящей
              подалгебре Ли\sourcefootnote{humphreys-ru,Kurosh1973}.
        \item Более трудная теорема Адо--Ивасавы устанавливает, что конечномерные алгебры Ли
              могут быть получены таким способом из конечномерных ассоциативных
              алгебр\sourcefootnote{jacobson-ru}.
        \item В случае, когда алгебра Ли $\mathfrak{L}(\Phi,\K)$ является простой,
              присоединенное представление точно, и ассоциативная алгебра, порожденная
              образом $N\Phi(\K)$, дает ее конечномерную обертывающую алгебру.
    \end{itemize}
\end{frame}

\begin{frame}[t]{Структурные константы в базисе Шевалле}
    \small
    Для нильтреугольных алгебр $N\Phi(\K)$ В.\,М.~Левчук построил семейства точных
    обертывающих алгебр\sourcefootnote{DAN2018}.

    \medskip

    По теореме Шевалле о базе при $r,s\in\Phi$ имеем\sourcefootnote{carter72}
    \[
        [e_r,e_s]=N_{rs}e_{r+s}=-[e_s,e_r]\quad (r+s\in\Phi),
        \qquad [e_r,e_s]=0\quad (r+s\notin \Phi\cup\{0\}),
    \]
    где $N_{rs}=\pm 1$ для типа $A_n$.

    \begin{itemize}
        \item Знаки структурных констант $N_{rs}$ имеют определенный произвол для алгебры
              Шевалле, а также для $N\Phi(\K)$.
    \end{itemize}
\end{frame}

\begin{frame}[t]{Семейство алгебр $R_\Phi$}
    \small
    \begin{block}{Определение}
        Фиксируя выбор знаков структурных констант $N_{rs}$, через $R_\Phi$ обозначим $\K$-алгебру с базой $\{e_r\mid r\in\Phi^+\}$ и умножением
        \[
            e_re_s=0 \quad (r+s\notin\Phi^+),
        \]
        \[
            e_r e_s=e_{r+s},\qquad e_s e_r=(1-N_{rs})e_{r+s}
            \quad (r,s,r+s\in\Phi^+,\ N_{rs}\ge 1).
        \]
    \end{block}

    \begin{itemize}
        \item Любую из построенных точных обертывающих алгебр $R_\Phi$ однозначно определяет
              выбор знака $\pm$ для каждой экстраспециальной пары корней.
    \end{itemize}
\end{frame}

\begin{frame}[t]{Тип $A_{n}$ и постановка задачи}
    \small
    \begin{block}{Естественный пример}
        Алгебра Ли $N\Phi(\K)$ типа $A_{n}$ изоморфна алгебре Ли,
        ассоциированной с алгеброй $NT(n+1,\K)$ нижних нильтреугольных матриц.
        Таким образом, для алгебры Ли $N\Phi(\K)$ типа $A_{n}$ существует
        ассоциативная точная обертывающая алгебра --- алгебра $NT(n+1,\K)$.
    \end{block}

    \begin{itemize}
        \item Ранее единственность $NT(n+1,\K)$ внутри семейства $R_\Phi$ исследовалась при
              дополнительном условии на идеалы и при предположении
              ассоциативности\sourcefootnote{imm2020,smj2023}.
        \item Мы усиливаем эти результаты, давая условие единственности ассоциативной точной
              обертывающей алгебры $N\Phi(\K)$ типа $A_n$: существование градуировки,
              согласованной с градуировкой Картана.
    \end{itemize}
\end{frame}

\begin{frame}[plain]
    \vfill
    \begin{center}
        {\usebeamerfont{title}\usebeamercolor[fg]{title}Градуированные алгебры и\\теорема единственности для типа $A_n$.\par}
    \end{center}
    \vfill
\end{frame}

\begin{frame}[t]{Градуированные алгебры}
    \small
    Пусть $G$ --- абелева группа. Будем говорить, что на алгебре $A$ задана
    \term{$G$-градуировка} $\Gamma$, если $A$ представлена в виде прямой суммы
    подпространств
    \[
        \Gamma\colon A=\bigoplus_{g\in G}A_g,
        \qquad A_gA_h\subset A_{g+h}.
    \]

    \medskip

    Подалгебра $R\subset A$ называется \term{градуированной}, если
    $R=\bigoplus_{g\in G}(A_g\cap R)$.

    \medskip

    Если $I$ является градуированным идеалом алгебры $A$, то фактор-алгебра $A/I$
    наследует градуировку\sourcefootnote{moody95}, задаваемую формулой
    \[
        (A/I)_g\coloneq (A_g+I)/I = \{a+I \mid a\in A_g\}.
    \]
\end{frame}

\begin{frame}[t]{Тонкие градуировки}
    \begin{itemize}
        \item Будем говорить, что $\Gamma_2 \colon A = \bigoplus_{h \in H}A_h$ является
              \emph{уточнением} градуировки $\Gamma_1 \colon A = \bigoplus_{g \in G}A_g$,
              если существует гомоморфизм огрубления $\phi\colon H\to G$ такой, что для
              каждого $h\in H$ выполняется включение $A_h\subset A_{\phi(h)}$.

        \item Если для некоторого $h\in H$ это включение строгое, то будем говорить о
              \emph{собственном уточнении}.

        \item Градуировка $\Gamma$ называется \term{тонкой}, если она не допускает
              собственных уточнений\sourcefootnote{patera1989,kochetov2013}.

        \item Тонкую градуировку $\Gamma$ будем называть \term{градуировкой ширины один},
              если все ее нетривиальные компоненты одномерны\sourcefootnote{mill-graded2019}.
    \end{itemize}
    Любая градуировка конечномерной алгебры допускает тонкое уточнение.
\end{frame}

\begin{frame}[t]{Градуировка Картана}
    \small
    Базис Шевалле алгебры $\mathfrak{L}(\Phi,\K)$ ранга $n$ задает на ней
    \term{$\mathbb Z\Phi$-градуировку
        Картана}\sourcefootnote{kochetov2013,moody95} целочисленной решеткой корней
    $\mathbb Z\Phi\simeq \mathbb Z^n$.

    \begin{itemize}
        \item Компонент с нулевой степенью в этой градуировке является подалгеброй Картана
              $H$.
        \item Остальные нетривиальные компоненты являются одномерными $H$-инвариантными
              подалгебрами в $\mathfrak L(\Phi,\K)$.
        \item В частности, на подалгебре $N\Phi(\K)$ индуцированная градуировка имеет вид
              \[
                  \Gamma_C\colon N\Phi(\K)=\bigoplus_{r\in\Phi^+}\K e_r,
              \]
              причем все ее нетривиальные компоненты одномерны.
    \end{itemize}
\end{frame}

\begin{frame}[t]{Эквивалентность градуировок}
    Рассмотрим две градуированные алгебры
    \begin{equation}\label{eq:two-graded-algebras}
        \Gamma_1 \colon A = \bigoplus_{g \in G}A_g
        \quad \text{и} \quad
        \Gamma_2 \colon B = \bigoplus_{h \in H}B_h.
    \end{equation}
    \begin{block}{Эквивалентность}
        Градуировки $\Gamma_1$ и $\Gamma_2$ называются \term{эквивалентными}, если существует изоморфизм алгебр $\phi\colon A\to B$ и для любого $g$ найдется $h$ такой, что $\phi(A_g)=B_h$.
    \end{block}
\end{frame}

\begin{frame}[t]{Тонкие градуировки нильпотентных алгебр}
    \begin{block}{Предложение}
        Любая тонкая градуировка конечномерной нильпотентной алгебры $A$ является градуировкой ширины один.
    \end{block}

    Доказательство предложения проводится индукцией по $\dim A$ с использованием
    лемм:

    \begin{itemize}
        \item Аннулятор $\Ann(A)$ градуированной алгебры является градуированным идеалом.

        \item Если на алгебре $A$ задана тонкая градуировка, то индуцированная градуировка на
              $\quot{A}{\Ann(A)}$ также тонкая.

        \item Если $A_g\cap \Ann(A)\ne 0$, то $\dim A_g=1$ и $A_g\subset\Ann(A)$.
    \end{itemize}
\end{frame}

\begin{frame}[t]{Лемма о восстановлении корневой структуры}
    На алгебре $N\Phi(\mathbb{K})$ зафиксируем градуировку Картана
    \[
        \Gamma_C\colon N\Phi(\mathbb{K})=
        \bigoplus_{r\in\Phi^+}\mathbb{K}e_r.
    \]
    \begin{block}{Лемма}
        Пусть $R$ --- нильпотентная точная обертывающая алгебра алгебры Ли $N\Phi(\K)$. Если на $R$ задана тонкая градуировка $\Gamma$ такая, что перенесенная на $R^{(-)}$ градуировка эквивалентна $\Gamma_C$, то $R$ можно реализовать на линейном пространстве $N\Phi(\K)$ так, что
        \[
            e_r e_s=0\quad \text{при } r+s\notin\Phi^+,
        \]
        а коммутатор в $R^{(-)}$ совпадает с коммутатором в $N\Phi(\K)$.
    \end{block}
\end{frame}

\begin{frame}[t]{Нильпотентность ассоциативных точных обертывающих алгебр}
    \begin{block}{Лемма}
        Пусть $L$ --- конечномерная нильпотентная алгебра Ли такая, что
        \[
            \dim Z(L)=1,\qquad Z(L)\subset [L,L].
        \]
        Тогда всякая ассоциативная точная обертывающая алгебра алгебры Ли $L$ является
        нильпотентной.
    \end{block}

    \begin{block}{Следствие}
        Пусть $R$ --- ассоциативная точная обертывающая $\K$-алгебра для алгебры Ли $N\Phi(\K)$ типа $A_n$ $(n\ge 2)$. Тогда $R$ нильпотентна.
    \end{block}
\end{frame}

\begin{frame}[t]{Основной результат для типа $A_n$}
    \begin{alertblock}{Теорема}
        Если ассоциативная точная обертывающая алгебра $R$ алгебры Ли $N\Phi(\K)$ типа $A_n$ $(n\ge 3)$ допускает градуировку, согласованную с градуировкой Картана, то
        \[
            R \simeq NT(n+1,\K).
        \]
    \end{alertblock}

    \begin{itemize}
        \item Теорема формулируется для произвольной ассоциативной точной обертывающей
              алгебры, а не только для специального семейства $R_\Phi$.
        \item При этом условие согласованности с $\Gamma_C$ заменяет дополнительные
              ограничения из более ранних результатов\sourcefootnote{imm2020,smj2023}.
    \end{itemize}
\end{frame}

\begin{frame}[t]{Существенность требования согласованности с $\Gamma_C$}
    \small
    Если в теореме для $A_n$ не требовать существования градуировки, согласованной с градуировкой Картана, то алгебра $R$, вообще говоря, не обязана быть изоморфной $NT(n+1,\K)$.

    \medskip

    В алгебре $NT(4,\K)$ заменим умножение на новое $\ast$, полагая $x\ast y=xy$
    для всех пар базисных элементов, кроме пары $(e_{32},e_{21})$, и дополнительно
    задавая
    \[
        e_{32}\ast e_{21}=e_{31}+e_{41},\qquad e_{21}\ast e_{32}=e_{41}.
    \]
    Тогда
    \begin{itemize}
        \item все коммутаторы базисных элементов совпадают с коммутаторами в $NT(4,\K)$;
        \item ассоциативность сохраняется, поскольку добавка лежит в аннуляторе;
        \item новая алгебра не изоморфна $NT(4,\K)$, поскольку меняется размерность левого
              аннулятора.
    \end{itemize}
\end{frame}

\begin{frame}[plain]
    \vfill
    \begin{center}
        {\usebeamerfont{title}\usebeamercolor[fg]{title}Левосимметрические структуры\par}
    \end{center}
    \vfill
\end{frame}

\begin{frame}[t]{Левосимметрические алгебры}
    \small
    Через $L_x\colon y\mapsto xy$ $(x,y \in R)$ обозначим линейные операторы левого
    умножения в алгебре $R$.

    \begin{block}{Определение}
        Алгебра $R$ называется \term{левосимметрической}, если
        \[
            (xy)z-x(yz)=(yx)z-y(xz)\qquad (x,y,z\in R),
        \]
        или, эквивалентно,
        \[
            [L_x,L_y]=L_{[x,y]}.
        \]
    \end{block}

    \begin{itemize}
        \item Как и в ассоциативном случае, всякая левосимметрическая алгебра Ли-допустима.
        \item Левосимметрическую точную обертывающую алгебру для алгебры Ли $\mathfrak L$
              называют \term{левосимметрической структурой} на $\mathfrak L$.
    \end{itemize}
\end{frame}

\begin{frame}[t]{Вопрос о существовании левосимметрических структур}
    \begin{itemize}
        \item Полупростая алгебра Ли над полем характеристики $0$ не допускает
              левосимметрических структур\sourcefootnote{burde-pre-lie-survey}.
        \item Дж.~Милнор поставил вопрос, эквивалентный вопросу о существовании
              левосимметрических структур для разрешимых алгебр
              Ли\sourcefootnote{milnor-left-sym}.
        \item Существуют нильпотентные алгебры Ли, не допускающие левосимметрических
              структур\sourcefootnote{burde-counter-example}.
    \end{itemize}
\end{frame}

\begin{frame}[t]{Полупрямые суммы алгебр Ли}
    Для построения левосимметрических точных обертывающих алгебр для $N\Phi(\K)$
    типов $B_n$, $C_n$ и $D_n$ используем разложения этих алгебр в полупрямые
    суммы подалгебры $NA_{n-1}(\K)$ с подходящим идеалом.

    \medskip

    \begin{block}{Определение\sourcefootnote{vinberg-onishik}}
        \term{Полупрямой суммой $\mathfrak g\ltimes_{\phi}\mathfrak h$} алгебр Ли $\mathfrak g$ и $\mathfrak h$ называется прямая сумма векторных пространств $\mathfrak g$ и $\mathfrak h$, снабженная операцией коммутирования
        \[
            [(x,a),(y,b)] = \bigl([x,y],\ \phi_x(b)-\phi_y(a)+[a,b]\bigr),
        \]
        где $\phi\colon\mathfrak{g}\to\Der(\mathfrak{h})$ --- некоторый гомоморфизм
        алгебры Ли $\mathfrak{g}$ в алгебру Ли дифференцирований $\Der(\mathfrak{h})$ и
        $\phi_x = \phi(x)$.
    \end{block}

\end{frame}

\begin{frame}[t]{Полупрямые суммы левосимметрических алгебр}
    \small
    \begin{block}{Определение}
        \term{Полупрямой суммой} $R\ltimes_\tau I$ левосимметрических алгебр $R$ и
        $I$ назовем прямую сумму векторных пространств $R$ и $I$, снабженную
        операцией умножения
        \[
            (x,a)(y,b)=(xy,\ \tau_x(b)+ab),
        \]
        где $\tau\colon R^{(-)} \to \Der(I)$ --- некоторый гомоморфизм алгебры Ли
        $R^{(-)}$ в алгебру Ли дифференцирований $\Der(I)$ и $\tau_x = \tau(x)$.
    \end{block}

    \medskip

    Всякое дифференцирование левосимметрической алгебры $I$ является
    дифференцированием алгебры Ли $I^{(-)}$; поэтому вместе с $R\ltimes_\tau I$
    можно одновременно рассматривать и полупрямую сумму алгебр Ли
    $R^{(-)}\ltimes_\tau I^{(-)}$.
\end{frame}

\begin{frame}[t]{Два результата о полупрямых суммах левосимметрических алгебр}
    \small
    \begin{block}{Теорема}
        %
        Пусть $R$ и $I$ --- левосимметрические алгебры и $\tau\colon R^{(-)}\to\Der(I)$
        --- гомоморфизм алгебры Ли $R^{(-)}$ в $\Der(I)$. Тогда алгебра
        $R\ltimes_{\tau} I$ является левосимметрической. Более того, $(R\ltimes_{\tau}
        I)^{(-)}$ изоморфна полупрямой сумме $R^{(-)}\ltimes_\tau I^{(-)}$.
    \end{block}

    \begin{block}{Следствие}
        %
        Пусть $\mathfrak{L} = \mathfrak{l}\ltimes_{\phi}\mathfrak{h}$ --- полупрямая
        сумма алгебр Ли, $R$ и $I$ --- левосимметрические точные обертывающие алгебры
        для $\mathfrak{l}$ и $\mathfrak{h}$ соответственно. Пусть также задан
        гомоморфизм алгебр Ли $\tau\colon R^{(-)}\to\Der(I)$ такой, что индуцированное
        действием $\tau$ действие алгебры Ли $R^{(-)}$ на $I^{(-)}$ совпадает с
        действием $\phi$. Тогда алгебра $R\ltimes_{\tau} I$ является левосимметрической
        точной обертывающей алгеброй алгебры Ли $\mathfrak{L}$.
    \end{block}
\end{frame}

\begin{frame}[t]{Когда полупрямая сумма ассоциативных алгебр ассоциативна?}
    \small
    \begin{block}{Теорема}
        Пусть $R$ и $I$ --- ассоциативные алгебры и $\tau\colon R^{(-)}\to\Der(I)$
        --- гомоморфизм алгебры Ли $R^{(-)}$ в $\Der(I)$. Алгебра
        $R\ltimes_{\tau} I$ ассоциативна тогда и только тогда, когда выполнены условия:
        \begin{enumerate}
            \item[\textup{(A)}] $\tau_{xy}=\tau_x\tau_y$ для всех $x,y\in R$;
            \item[\textup{(B)}] $\tau_x(I)\subset \Ann_r(I)$ для всех $x\in R$.
        \end{enumerate}
    \end{block}
\end{frame}

\begin{frame}[t]{Представления корней типов $B_n$, $C_n$, $D_n$}
    \small
    %
    Пусть $\Psi=A_{n-1}^+=\{\epsilon_i-\epsilon_j \mid 1\leq j<i\leq n\}$.
    Тогда\sourcefootnote{bourbaki72}
    \[
        D_n^+=\Psi\cup\{\epsilon_i+\epsilon_j \mid 1\le j<i\le n\},
    \]
    \[
        C_n^+=D_n^+\cup\{2\epsilon_i \mid 1\le i\le n\},
        \qquad
        B_n^+=D_n^+\cup\{\epsilon_i \mid 1\le i\le n\}.
    \]

    \begin{itemize}
        % \item Таким образом, $\Phi^+$ состоит из корней вида
        %       $r_{i,mj}=\epsilon_i-m\epsilon_j$, где $m\in\{0,-1,+1\}$.
        \item Базу $\Pi(\Phi)$ запишем в виде $\Pi(\Psi)\cup\{r_\Phi\}$, где
              \[
                  r_\Phi=
                  \begin{cases}
                      \epsilon_2 + \epsilon_1, & \Phi=D_n, \\
                      2\epsilon_1,             & \Phi=C_n, \\
                      \epsilon_1,              & \Phi=B_n.
                  \end{cases}
              \]
        \item Подалгебра $N\Psi$ изоморфна $NT(n,\K)^{(-)}$.
    \end{itemize}
\end{frame}

\begin{frame}[t]{Разложение $N\Phi(\K)$ типов $B_n$, $C_n$ и $D_n$ в полупрямую сумму}
    \small
    Для корней будем писать $s\ge r$, если все коэффициенты в разложении $s-r$ по базе простых корней $\Pi$ неотрицательны. Определим идеалы
    \[
        T(r)\coloneq \sum_{s\ge r}\K e_s.
    \]

    \begin{itemize}
        \item Обозначим через $\phi$ присоединенное действие алгебры Ли $NT(n,\K)^{(-)}$ на
              идеале $T(r_\Phi)$.
        \item Тогда для типов $B_n$, $C_n$ и $D_n$ получаем разложение алгебр Ли
              \[
                  N\Phi(\K)\simeq NT(n,\K)^{(-)}\ltimes_\phi T(r_\Phi).
              \]
    \end{itemize}
\end{frame}

\begin{frame}[t]{Выбор умножения на идеале $T(r_\Phi)$}
    \small
    \begin{center}
        \begin{tabular}{>{\raggedright\arraybackslash}p{0.12\linewidth} >{\raggedright\arraybackslash}p{0.27\linewidth} >{\raggedright\arraybackslash}p{0.51\linewidth}}
            \toprule
            Тип   & Идеал                        & Выбор ассоциативного умножения                  \\
            \midrule
            $D_n$ & $T(\epsilon_2 + \epsilon_1)$ & идеал абелев; берем нулевое умножение           \\
            $C_n$ & $T(2\epsilon_1)$             & идеал абелев; берем нулевое умножение           \\
            $B_n$ & $T(\epsilon_1)$              & при $\Char\K\neq 2$ полагаем
            $\displaystyle e_r\cdot e_s=\frac12[e_r,e_s]$; при $\Char\K=2$ берем нулевое умножение \\
            \bottomrule
        \end{tabular}
    \end{center}

    \begin{itemize}
        \item Для типа $B_n$ используется то, что коммутант алгебры Ли $T(\epsilon_1)$ лежит
              в ее центре, поэтому
              \[
                  T(\epsilon_1)^3=0.
              \]
        \item Во всех трех случаях $(T(r_\Phi),\cdot_\Phi)^{(-)}\simeq T(r_\Phi)$, а действие
              $\phi_x$ задает дифференцирование этой ассоциативной алгебры.
    \end{itemize}
\end{frame}

\begin{frame}[t]{Главный результат для типов $B_n$, $C_n$, $D_n$}
    \begin{alertblock}{Теорема}
        Полупрямая сумма ассоциативных алгебр
        \[
            NT(n,\K)\ltimes_\phi T(r_\Phi)
        \]
        является левосимметрической точной обертывающей алгеброй для $N\Phi(\K)$ типов
        $B_n$, $C_n$ и $D_n$.
    \end{alertblock}
\end{frame}

\begin{frame}[t]{Почему эти конструкции неассоциативны}
    \small
    \begin{block}{Тип $B_n$}
        При $\Char\K\neq 2$ и $n\ge 3$ выполнение условия
        \[
            \phi_x\bigl(T(\epsilon_1)\bigr)\subset \Ann_r\bigl(T(\epsilon_1),\cdot\bigr)
        \]
        невозможно: образ действия $NT(n,\mathbb{K})^{(-)}$ на факторе
        $\quot{T(\epsilon_1)}{Z}$ имеет размерность не менее $2$, а проекция правого
        аннулятора в $\quot{T(\epsilon_1)}{Z}$ --- не более $1$.
    \end{block}

    \begin{block}{Типы $C_n$ и $D_n$}
        %
        Здесь нарушается условие \[\phi_{xy}=\phi_x\phi_y.\] При $x=e_{\epsilon_2 - \epsilon_1}$, $y=e_{\epsilon_3 - \epsilon_2}$ и
        $u=e_{\epsilon_2 + \epsilon_1}\in T(r_\Phi)$ имеем $\phi_{xy}(u)=0$, но $\phi_x\phi_y(u)\ne 0$.
    \end{block}
\end{frame}

\begin{frame}[t]{Итоги}
    \begin{itemize}
        \item Для типа $A_n$ ассоциативная точная обертывающая алгебра, допускающая
              градуировку, согласованную с $\Gamma_C$, единственна с точностью до
              изоморфизма:
              \[
                  R\simeq NT(n+1,\K).
              \]
        \item Для типов $B_n$, $C_n$ и $D_n$ построены градуированные левосимметрические
              точные обертывающие алгебры в виде полупрямых сумм.
        \item В этих конструкциях ассоциативность, вообще говоря, теряется: в типе $B_n$ при
              $\Char\K\neq 2$ и $n \geq 3$, в типе $C_n$ при $n\ge 3$, в типе $D_n$ при $n\ge
              4$.
    \end{itemize}
\end{frame}

\end{document}
